% Chapter 2: Literature Review
% Covers thesis.md §4 (Literature Survey: Existing Algorithms)

This chapter reviews the existing algorithms for mixing-tree construction on digital microfluidic biochips. The field spans over 15 years of research. Only RMA and BS are implemented in this work as baselines; the remaining algorithms are discussed as related work.

\section{RMA --- Ratioed Mixing Algorithm}
\label{sec:rma}

\textbf{Reference:} Roy, Bhattacharya, Chakrabarti, and Chakrabarty~\cite{roy2010layout} proposed the Ratioed Mixing Algorithm (RMA) for layout-aware solution preparation on DMFBs.

\textbf{Strategy:} RMA is a top-down greedy algorithm. At each node, the \textbf{dominant fluid} (highest count) is assigned first to one child up to volume $L/2$; remaining capacity is filled by splitting the next-largest fluid as needed.

\textbf{Strength:} Long dilution chains (high $l$ and $\text{maxL}$); good for hardware layouts where serial dilution is cheap.

\textbf{Weakness:} Greedy myopia---a locally optimal dominant-fluid pick can force more splits or mixes downstream than necessary.

\section{BS --- Bit-Scanning}
\label{sec:bs}

\textbf{Reference:} Thies et al.~\cite{thies_bs} proposed the Bit-Scanning (BS) approach for mixture generation on DMFBs.

\textbf{Strategy:} BS is a bottom-up, binary-encoding-driven algorithm. For each fluid $f$ with count $c_f$, the bit-pattern of $c_f$ indicates at which tree levels a leaf for $f$ must appear. Construction proceeds level-by-level, pairing leaves and previously-built subtrees.

\textbf{Strength:} Minimum leaf count and fewest mixes among the baselines (every count is encoded in $O(\log c_f)$ leaves).

\textbf{Weakness:} No awareness of fluid splits or dilution structure; trees have higher fluid scattering.

\section{REMIA --- Reactant Minimization Algorithm}

\textbf{Reference:} Hsieh et al.~\cite{hsieh_remia} proposed REMIA, which uses a skewed mixing tree combined with an exponential dilution tree, designed specifically to \textbf{minimize reactant usage}.

\textbf{Trade-off:} REMIA often produces more waste droplets than IDMA/DMRW even though it consumes fewer total reagents.

\textbf{Status:} Not implemented in this work. Listed as primary related work; recommended as a future-work baseline.

\section{WARA --- Waste Recycling Algorithm}

WARA addresses multi-target sample preparation. It builds REMIA-style mixing trees per target, then identifies opportunities to recycle waste droplets from one target's tree as inputs to another.

\textbf{Status:} Not implemented; multi-target sample preparation is out of scope for this work, which focuses on single-target preparation.

\section{IDMA / DMRW}

The Improved Dilution and Mixing Algorithm (IDMA) and Dilution and Mixing with Reduced Wastage (DMRW)~\cite{waste_aware_2014} both target the dilution subproblem (one reagent plus one buffer) by carefully tracking waste-droplet recycling within a single tree.

\textbf{Status:} Not implemented; their domain (single-reagent dilution) is a special case of the general multi-fluid mixing addressed here.

\section{MinMix and CoDOS}

MinMix and CoDOS are generic mixing-tree construction baselines from earlier work, frequently cited alongside RMA.

\textbf{Status:} Not implemented; RMA is used as the representative greedy baseline.

\section{Skewed Mixing Trees}

Mitra et al.~\cite{mitra2013skewed} proposed using deliberately unbalanced (skewed) trees to exploit reagent-asymmetry and reduce reactant volume.

\textbf{Status:} Not implemented; orthogonal in tree-shape philosophy to this work.

\section{Layout-Aware Mixture Preparation}

Bhattacharjee et al.~\cite{bhattacharjee2015layout} coupled mixing-tree construction to physical chip layout, jointly optimizing routing distance and reagent count.

\textbf{Status:} Not implemented; this work is layout-agnostic at present (a clear future-work direction).

\section{Demand-Driven Preparation with Droplet Streaming}

This approach~\cite{dac2014demand} optimizes for repeated mixture demand (e.g., master-mix in PCR) rather than single-target output.

\textbf{Status:} Not implemented.

\section{ILP-based Synthesis for Sample Preparation}

The ILP-based synthesis approach~\cite{ilp_ieee2016} models the sample-preparation problem as an integer linear program and solves it with a commercial solver (CPLEX/Gurobi).

\textbf{Relevance to this work:} This is the \textit{solver-based} ILP approach, \textbf{distinct from our ``ILP'' algorithm} described in Chapter~3, which is an integer-allocation \textbf{DP heuristic} (no LP relaxation, no commercial solver). We share the integer-allocation framing but not the solution method. We do not directly compare against this work because we do not have the solver setup.

\section{SIMOP --- Simulation-Guided Optimization}

SIMOP~\cite{simop2021} couples tree generation with simulated chip behavior to optimize accuracy and error management from unbalanced splits.

\textbf{Status:} Not implemented.

\section{Summary of Literature Coverage}

Table~\ref{tab:lit_coverage} summarizes the algorithm families in the literature and their implementation status in this work.

\begin{table}[H]
\centering
\caption{Summary of algorithm families and implementation status.}
\label{tab:lit_coverage}
\begin{tabular}{|l|l|l|}
\hline
\textbf{Family} & \textbf{Algorithms} & \textbf{Implemented?} \\
\hline
Greedy top-down & RMA, MinMix, CoDOS & RMA only \\
\hline
Bottom-up bitwise & BS & Yes \\
\hline
Reactant-minimal & REMIA, IDMA, DMRW & No (future work) \\
\hline
Multi-target & WARA & No (out of scope) \\
\hline
Layout-aware & Bhattacharjee et al. & No (future work) \\
\hline
Skewed-tree & Mitra et al. & No \\
\hline
Solver-based & ILP-based Synthesis & No (heuristic ILP only) \\
\hline
Simulation-guided & SIMOP & No \\
\hline
\end{tabular}
\end{table}

The two implemented baselines---RMA and BS---represent the two extremes of the splits-versus-dilution trade-off: RMA optimizes for long dilution chains at the cost of more splits, while BS minimizes mixes and splits at the cost of destroyed dilution structure. The proposed algorithms in Chapter~3 aim to fill the gap between these extremes.
